% ---
% id: reductions/def:emitnorm
% kind: def
% title: "The emission normalizer, the declared datum and the declared model"
% summary: >-
%   The normalizer puts an exact datum back on the lattice class by class: it
%   keeps the classes already there and snaps the others for fidelity, rounding
%   to nearest and placing the residue only where support allows, or, failing
%   that, by a flagged last resort.
% prominence: standard
% status: implemented
% uses: [lem:exactpsd, def:exactadm, def:reduction-instance, def:joint-family, foundations/def:snap, foundations/def:exact-reference]
% mentions: [thm:emit, thm:instrument-instance, cor:reference-companion, foundations/req:exact-emission, foundations/req:who-snaps, foundations/req:snap-reference, foundations/req:support-safe-snap, foundations/ass:snapped-measurement, req:emission-record, req:emission-refusal, meat-core/frank-wolfe/req:certified-evaluation]
% assumes: []
% refs: []
% provenance: newly-derived
% implements: certiqs_key.reductions.normalizer:normalize
% lean:
% notes: "Two readings of foundations/req:support-safe-snap (iv) are used by the last resort (3): the family mixed is the class of the handed datum before any step (the exact reference), and the uniform family is a positive multiple of the identity on the elements in the sector's support only, the elements outside it staying exactly zero. Step (2) mixes the handed datum itself with its diagonal parts; that is an instance of foundations/req:support-safe-snap (ii) under its reading that admits the data D themselves where every entry is a Gaussian rational, which every entry of the handed datum is. The procedure aims at fidelity and at nothing else: no claim relates a snapped class to the handed datum, and a run evaluates off the snap through the reference companion (cor:reference-companion), which bypasses the normalizer."
% ---

A \tref{reduction instance} forms its \tref{handed datum} by exact
arithmetic; every exact step multiplies, and a product of dyadics adds
significand widths, so the datum can leave the declared lattice on which this
module emits. The normalizer defined here puts it back by a snap of the
module's own emission, never of anything it was handed. Every class of the
handed datum is Hermitian, positive semidefinite and complete at its class
weight with Gaussian-rational entries, so it is an \tref{exact reference} of
itself (\xref{foundations/def:exact-reference}, every radius zero). The
normalizer keeps the classes already on the lattice and snaps every other
class by one procedure that stays as close to the handed datum as it can
while keeping every row and column that the handed datum holds at zero
(\xref{foundations/req:support-safe-snap}): it rounds to nearest, softens the
off-diagonal entries where positivity needs it, and places what completeness
leaves only on elements whose support holds it. Where no such placement
exists, a last resort mixes the class with a multiple of the identity
(\xref{foundations/req:support-safe-snap}\,(iv)); that may give an element a
row the handed datum denies it, so the class is flagged. No claim relates a
snapped class to the handed datum. A run that forms the number its key
length is formed from on the reference companion
(Corollary~\ref{cor:reference-companion}), which the reductions build with no
normalizer from the \tref{handed reference}, or from the instance's own
snapped blocks where none is handed, evaluates every class off its snap; a run that does not, such as one whose
key length is formed from a number certified on the snapped declaration in
the fallback of \xref{meat-core/frank-wolfe/req:certified-evaluation}\,(i),
rests, on the classes the
normalizer moved, on \xref{foundations/ass:snapped-measurement}, and a run
that declares the strict option refuses there instead
(\xref{foundations/req:snap-reference}\,(iv)). The emission theorem
(Theorem~\ref{thm:emit}) is stated about what the normalizer produces. The
six conditions its last step decides are those of
\xref{foundations/req:exact-emission}; when a producer re-establishes its
lattice at an emission, and on what it may snap at all, is
\xref{foundations/req:who-snaps}. Everything the normalizer records is
Requirement~\ref{req:emission-record}, and its refusals are
Requirement~\ref{req:emission-refusal}.

\begin{definition}
\label{def:emitnorm}
Let $\mathcal{J}$ be a \tref{reduction instance}
(Definition~\ref{def:reduction-instance}) with \tref{handed datum}
$A\in\mathcal{D}_{\mathrm{adm}}$, formed exactly. Those numbers \emph{are} the
\tref[data sector]{data sectors} of the emitted target
(Definition~\ref{def:reduction-instance}\,(I2)), so they are this module's own
declaration to its \tref{consumer}. Where the exact arithmetic that formed
them has left the declared lattice \eqref{eq:latinv}, they
are re-established on it, per data sector $s$ and per class $t$, by the
following steps, the
\term[The deterministic procedure that puts a reduction instance's exactly formed datum back on the declared lattice: it keeps the classes already there, snaps every other class for fidelity without reviving a zero row or column or, failing that, by a flagged last resort, and decides the six emission conditions last.]{emission normalizer}.
The steps act on the elements $A^{(s)}_v$ with
$v\in\Vset_t\cap\Vset^{(s)}$, the elements of the class $t$ in the sector's
support; the elements \tref{outside the sector support} are exactly zero,
lie on the lattice, and are touched by no step. Every positive
semidefiniteness below is decided exactly by \eqref{eq:psdtest}, on matrices
with entries in $\Fldi$. The
\tref{structural support} of a matrix is the set of indices $i$ whose row $i$
or column $i$ holds a non-zero entry (\xref{foundations/def:snap}). The
\emph{designation order} ranks the elements of a class in the sector's
support by largest exact trace and, among equal traces, by the
\tref{declared order} of the class. \emph{A tie is exact and ordinary here}:
the traces compared are exact rationals, and a sector symmetric under a
relabelling of the class ties every element of it, so ``largest trace''
alone names no element; the declared order is the order in which this
contract reads $\Vset$, the outcome varying fastest, and it is emitted with
the artefact, so the order is a function of the declaration and of nothing
else. \emph{Rounding to nearest} a matrix at $\kappa$ rounds every real and
imaginary part of every entry to the nearest integer multiple of
$2^{-\kappa}$, \emph{ties to even}: a part lying exactly halfway between two
multiples goes to the one whose integer numerator is even. \emph{The tie is
exact}: the rounding is carried out in the substrate of
Lemma~\ref{lem:exactpsd} on exact rationals, so a halfway part is an exact
equality rather than a floating-point near-miss. \emph{Ties to even is
odd-symmetric}: negating a part negates both of the multiples nearest it and
their numerators, which preserves evenness, so the rounding commutes with
$X\mapsto X^{\dagger}$ --- mirror entries have equal real parts and opposite
imaginary parts before it and therefore after it --- and a matrix rounded
from a Hermitian one is Hermitian; a zero entry stays zero.
\begin{enumerate}[label=(\arabic*),leftmargin=2.4em]
\item \emph{Keep.} A class $t$ at a sector $s$ every element of which
  satisfies \eqref{eq:latinv} entry by entry is \emph{kept}:
  $A'^{(s)}_v\coloneqq A^{(s)}_v$ for every $v\in\Vset_t\cap\Vset^{(s)}$, and
  no further step acts on it.
\item \emph{Snap for fidelity.} At every other class and sector, for $w=0$
  and then $w=2^{-\kappa},2^{-\kappa+1},\dots,2^{-1}$ in this order, let
  $P^{w}_v$ be $A^{(s)}_v$ with every off-diagonal entry multiplied by
  $1-w$, that is $(1-w)A^{(s)}_v+w\,T_v$ with $T_v$ the diagonal part of
  $A^{(s)}_v$. Every element of the class in the sector's support receives
  $P^{w}_v$ rounded to nearest at $\kappa$, and the \emph{class residue}
  $E^{w}\coloneqq q_t\idon{d_s}-\sum_{v\in\Vset_t\cap\Vset^{(s)}}A'^{(s)}_v$,
  a Hermitian matrix of \tref[lattice value]{lattice values}, is split by its \emph{components}:
  the classes of the equivalence relation on the indices $i$ whose row $i$ of
  $E^{w}$ is non-zero generated by $i\sim l$ whenever $E^{w}_{il}\ne0$. For a
  component $C$ write $E_C$ for the matrix that agrees with $E^{w}$ on the
  entries whose row and column both lie in $C$ and is zero elsewhere. Each
  $E_C$ is added to the first element $A'^{(s)}_u$ in the designation order
  whose structural support in the handed datum, that of $A^{(s)}_u$,
  contains $C$, its \emph{host}. The first $w$ at which every component has
  a host and every element so formed is decided positive semidefinite gives
  the family, and $w$ is recorded. Where no $w$ does, the class is snapped
  by step~(3).
\item \emph{The last resort.} The support-safe steps of (2) produced no snap
  of the class, and the class is snapped by a uniform lift, applied to the
  class with constant $q_t$ and starting from the handed datum. Let $n$ be
  the number of elements of the class in the sector's support, $v_1$ the
  first of them in the designation order, and, for $w\in(0,1)$,
  \begin{equation}\label{eq:lastresort}
    B^{w}_v\;\coloneqq\;(1-w)\,A^{(s)}_v+w\,\frac{q_t}{n}\,\idon{d_s},
    \qquad v\in\Vset_t\cap\Vset^{(s)},
  \end{equation}
  a family that is Hermitian, complete at $q_t$, with entries in $\Fldi$, and
  positive definite in every element, since $q_t>0$ at a class that is not
  kept (at $q_t=0$ completeness and positive semidefiniteness make every
  element zero, which is on the lattice). Its designation order is that of
  the class, $\tr B^{w}_v$ being an increasing affine function of
  $\tr A^{(s)}_v$. For $w=2^{-\kappa},2^{-\kappa+1},\dots,2^{-1}$ in this
  order, every element other than $v_1$ receives $B^{w}_v$ rounded to
  nearest at $\kappa$, and $v_1$ receives $q_t\idon{d_s}$ minus their sum;
  the first $w$ at which every element so formed is decided positive
  semidefinite gives the family. Where no $w$ does, every element $v\ne v_1$
  receives the first matrix $A'^{(s)}_v$ of the following list that is
  Hermitian, satisfies \eqref{eq:latinv} entry by entry, and satisfies
  \begin{equation}\label{eq:rounddown}
    0\;\preceq\;A'^{(s)}_v\;\preceq\;B^{1/2}_v :
  \end{equation}
  $B^{1/2}_v$ itself; then, for $r=\kappa,\kappa-1,\dots,1$ in this order,
  the shift $B^{1/2}_v-2^{-r}\idon{d_s}$ and the scaling
  $(1-2^{-r})B^{1/2}_v$, each rounded to nearest at $\kappa$; finally the
  zero matrix, which always qualifies, $B^{1/2}_v$ being positive
  semidefinite; and $v_1$ receives $q_t\idon{d_s}$ minus their sum. The
  weight $w$, or the fall-back, is recorded, and the class is flagged as one
  whose support is not preserved: an element may then hold a non-zero entry
  in a row or column in which the handed element holds none, that is,
  outside its structural support (\xref{foundations/def:snap}).
\item \emph{All six conditions are decided last}, after (2) or (3), each of
  which changes partial sums from the touched element onward: (a)--(e) of
  Definition~\ref{def:exactadm} together with the lattice invariant
  \eqref{eq:latinv} of every entry, read on the \tref{target family}
  \eqref{eq:instance-target} the datum enters. They are the six
  Theorem~\ref{thm:emit}\,(ii) decides; a family for which one of them fails
  is refused.
\end{enumerate}
A class at a sector that step~(1) does not keep is \emph{moved}, whichever
of (2) and (3) snapped it. Write $A'$ for the resulting datum --- the
\term[The datum the emission normalizer outputs: exactly admissible on the declared lattice, and carried entry by entry in the emitted target.]{declared datum},
an element of $\mathcal{D}_{\kappa}$ by (1) or the step of (2) and (3) that
snapped the class, each of which makes it Hermitian, on the lattice and
\tref{class complete} exactly, positive semidefinite by an exact decision in
(2) and in (3) at a weight of its ladder, and by construction in the fall-back of (3), and
by (4), which decides all six --- and
\begin{equation}\label{eq:declmodel}
  \Mdec\;\coloneqq\;M[A'],
  \qquad\text{that is}\qquad
  \Mdec_v\;=\;M_v[A']\quad(v\in\Vset),
\end{equation}
for the
\term[The handed model with its datum replaced by the declared datum, as the instance's model map defines it.]{declared model}:
the \tref{handed model} with the datum replaced by $A'$, as the instance's
\tref{model map} defines it. Per data sector $s$, those being the sectors the
normalizer touches and the only ones both data have,
\begin{equation}\label{eq:displacement}
  \Delta_s\;\coloneqq\;\max_{v\in\Vset}\
  \mxn{A'^{(s)}_v-A^{(s)}_v}
\end{equation}
is the
\term[Per data sector, the largest entrywise modulus of the difference between the declared datum and the handed datum.]{emission displacement};
it is recorded per data sector, together with the kept and the moved classes,
for a moved class the step that snapped it, the weight $w$ of (2) or (3) or
the fall-back of (3), the hosts of the components of (2), and the flag of a
class snapped by (3).

The normalizer has five properties, each of which Theorem~\ref{thm:emit}, or
a claim that instantiates it, uses.
\emph{First, $A'=A$ and every $\Delta_s=0$ whenever the exact construction
already satisfies \eqref{eq:latinv}}: every class is then kept, the
normalizer is the identity and nothing below has an antecedent. That is the
ordinary case at a nominal declaration and is not the general one: every
exact step of a construction multiplies, a product of dyadics adds
significand widths, and a lattice established at the instance is closed
under none of them; that is the case in which a producer re-establishes its
lattice at an emission.
\emph{Second, the normalizer reads data sectors and nothing else.} No object
the \tref{hypotheses (H)} of the instance are about --- an operator without a finite
presentation, a family on a part of $\Hs{\Bsys}$ the datum does not describe,
a register POVM --- is read, and none is repaired; so (H) holds of $\Mdec$
by the \tref{stability clause} of
Definition~\ref{def:reduction-instance}\,(I1), $\Mdec$ is a POVM on $\Bsys$
by (I1) at $A'$, and every statement the instance makes about its model under
(H) holds of $\Mdec$. For an instance whose datum is a family of preserved
blocks below a cutoff, this is the statement that the declared model agrees
with the handed one above the cutoff and satisfies every above-cutoff fact the
handed one does.
\emph{Third, what the normalizer produces is a datum, and the family the
emitted identity is about is the model the instance builds from it.} Two
objects, two symbols, and every statement below names which it is on: $A'$ is
the declared datum, a finite list of sector blocks carried entry by entry in
the emitted target; $\Mdec$ is the declared model on $\Bsys$, the object an
adjoint landing on $\Bsys$ can equal (Theorem~\ref{thm:emit}\,(iii)). Where
the model map is the identity on a finite $\Bsys$ the two coincide; where the
model has parts the datum does not describe --- an above-cutoff family,
further branches --- they do not, and the instance's model map says how $A'$
sits inside $\Mdec$. For a branching instance
(Theorem~\ref{thm:instrument-instance}) the normalizer acts branch by branch,
so the declared model is the branches' declared models composed.
\emph{Fourth, this is a snap of what this module declares and never of what it
was handed}, which is the condition under
which a producer may snap at all.
\emph{Fifth, no step revives a zero row or column, save the last resort
(3)}: for every data sector $s$ and every $v\in\Vset$ whose class at $s$ is
not snapped by (3) the structural support of $A'^{(s)}_v$ is contained in
that of $A^{(s)}_v$. In (2) the family $\{T_v\}$ of diagonal parts is
complete at $q_t$, the diagonal of $\sum_vA^{(s)}_v=q_t\idon{d_s}$ being
$q_t$ in every place, and $T_v$ is diagonal, with non-negative entries,
non-zero only at indices of the structural support of $A^{(s)}_v$, so
$P^{w}$ is a mixture with a complete diagonal family that adds only on
coordinates of each element's own structural support; $P^{w}_v$ has the
structural support of $A^{(s)}_v$, and a rounding to nearest keeps a zero
entry zero. Each $E_C$ has its non-zero rows and columns in $C$, which the
host's structural support contains, and the entries of $E^{w}$ between two
components are zero by their definition, so $E^{w}=\sum_CE_C$ and every
non-zero entry of the residue goes to a host. In (2) the residue need not be
positive semidefinite, and the family is taken only where every element is
decided positive semidefinite. No element is mixed with a multiple of the
identity outside (3). In (3) the elements of the class in the sector's
support may gain rows and columns, which is what its flag discloses; the
elements outside the support stay zero there too. And (3) always yields a
family whose elements are positive semidefinite: at a weight of its ladder by decision, and
in the fall-back because every element of $B^{1/2}$ is positive definite, the
rounded elements satisfy \eqref{eq:rounddown}, and
$A'^{(s)}_{v_1}=B^{1/2}_{v_1}+\sum_{v\ne v_1}\big(B^{1/2}_v-A'^{(s)}_v\big)
\succeq B^{1/2}_{v_1}\succ0$ by the completeness of $B^{1/2}$.
\end{definition}
